
\documentclass[10pt]{article}
\usepackage{lipsum}
\usepackage{booktabs}
\usepackage{subcaption}
\usepackage{tabularx}
\usepackage{amsmath, amssymb}
\usepackage{import}
\usepackage{bm}
\usepackage{listings}
\usepackage{geometry}
\usepackage{graphicx}
\usepackage[colorlinks=true,linkcolor=blue,urlcolor=blue,bookmarksnumbered=true]{hyperref}
\newcommand{\sourcecode}[2]{\href{https://github.com/tiagovla/fdtd.py/blob/7027c252cd23b6cfb8cb197e361dacfac9fd1c72/#1}{\texttt{#2}}}

\geometry{
	a4paper,
	total={170mm,257mm},
	left=20mm,
	top=20mm,
}

%%%Author definitions
\newcommand{\executeiffilenewer}[3]{%
	\ifnum\pdfstrcmp%
		{\pdffilemoddate{#1}}%
		{\pdffilemoddate{#2}}%
		>0%
			{\immediate\write18{#3}}%
	\fi%
}

\newcommand{\includesvg}[2][]{%
	\executeiffilenewer{#1#2.svg}{#1#2.pdf}%
	{inkscape -z -D --file=#1#2.svg --export-pdf=#1#2.pdf --export-latex}%
	\subimport{#1}{#2.pdf_tex}% s
}
\newcommand{\test}[2][]{
	% \immediate\write\{inkscpe}

}

\title{\textbf{Theory\\ Finite-Difference Time-Domain}}
\author{@tiagovla}
\date{}

\begin{document}
\maketitle
\section{Maxwell Equations}%
\subsection{Material and source conventions}
We assume linear, nondispersive media with time-independent, positive
permittivity and permeability and nonnegative conductivities. The vector
equations below use scalar parameters for isotropic media. The component
equations also apply to diagonal constitutive tensors aligned with the grid.
For each component $\alpha\in\{x,y,z\}$,
\begin{equation}
    \epsilon_\alpha=\epsilon_0\epsilon_{r,\alpha},
    \qquad \mu_\alpha=\mu_0\mu_{r,\alpha}.
\end{equation}
Thus $\epsilon_\alpha$ and $\mu_\alpha$ are absolute parameters in
$\mathrm{F/m}$ and $\mathrm{H/m}$. The impressed electric current density
$\vec{J}_i$ has units $\mathrm{A/m^2}$, and the equivalent impressed magnetic
current density $\vec{M}_i$ has units $\mathrm{V/m^2}$. Electric conductivity
$\sigma^e$ has units $\mathrm{S/m}$; the magnetic loss parameter $\sigma^m$
has units $\Omega/\mathrm{m}$. For ordinary media without equivalent magnetic
sources or magnetic loss, $\vec{M}_i=0$ and $\sigma^m=0$. Dispersion,
time-varying materials, and off-diagonal tensor coupling require additional
constitutive equations. The lumped-element currents are derived separately
in Section~\ref{sec:lumped}.

\subsection{Equations}%

\begin{equation}
	\frac{\partial \vec{E}}{\partial t} =
	\frac{1}{\epsilon} \left( \nabla \times \vec{H} - \sigma^{e} \vec{E} - \vec{J}_{i} \right)
\end{equation}

\begin{equation}
    \begin{aligned}
        \frac{\partial E_{x}}{\partial t} &= \frac{1}{\epsilon_{x}} \left( \left[ \nabla \times \vec{H} \right] \cdot \hat{x} - \sigma^{e}_{x} E_{x} - J_{ix} \right) \\
        & = \frac{1}{\epsilon_{x}} \left( \frac{\partial H_{z}}{\partial y} - \frac{\partial H_{y}}{\partial z} - \sigma^{e}_{x} E_{x} - J_{ix} \right)
    \end{aligned}
\end{equation}

\begin{equation}
    \begin{aligned}
        \frac{\partial E_{y}}{\partial t} &= \frac{1}{\epsilon_{y}} \left( \left[ \nabla \times \vec{H} \right] \cdot \hat{y} - \sigma^{e}_{y} E_{y} - J_{iy} \right) \\
        & = \frac{1}{\epsilon_{y}} \left( \frac{\partial H_{x}}{\partial z} - \frac{\partial H_{z}}{\partial x} - \sigma^{e}_{y} E_{y} - J_{iy} \right)
    \end{aligned}
\end{equation}

\begin{equation}
    \begin{aligned}
        \frac{\partial E_{z}}{\partial t} &= \frac{1}{\epsilon_{z}} \left( \left[ \nabla \times \vec{H} \right] \cdot \hat{z} - \sigma^{e}_{z} E_{z} - J_{iz} \right) \\
        & = \frac{1}{\epsilon_{z}} \left( \frac{\partial H_{y}}{\partial x} - \frac{\partial H_{x}}{\partial y} - \sigma^{e}_{z} E_{z} - J_{iz} \right)
    \end{aligned}
\end{equation}



\begin{equation}
	\frac{\partial \vec{H}}{\partial t} =
	\frac{1}{\mu} \left( - \nabla \times \vec{E} - \sigma^{m} \vec{H} - \vec{M}_{i} \right)
\end{equation}

\begin{equation}
    \begin{aligned}
        \frac{\partial H_{x}}{\partial t} &= \frac{1}{\mu_{x}} \left( - \left[ \nabla \times \vec{E} \right] \cdot \hat{x} - \sigma^{m}_{x} H_{x} - M_{ix} \right) \\
        & = \frac{1}{\mu_{x}} \left( \frac{\partial E_{y}}{\partial z} - \frac{\partial E_{z}}{\partial y} - \sigma^{m}_{x} H_{x} - M_{ix} \right)
    \end{aligned}
\end{equation}

\begin{equation}
    \begin{aligned}
        \frac{\partial H_{y}}{\partial t} &= \frac{1}{\mu_{y}} \left( - \left[ \nabla \times \vec{E} \right] \cdot \hat{y} - \sigma^{m}_{y} H_{y} - M_{iy} \right) \\
        & = \frac{1}{\mu_{y}} \left( \frac{\partial E_{z}}{\partial x} - \frac{\partial E_{x}}{\partial z} - \sigma^{m}_{y} H_{y} - M_{iy} \right)
    \end{aligned}
\end{equation}

\begin{equation}
    \begin{aligned}
        \frac{\partial H_{z}}{\partial t} &= \frac{1}{\mu_{z}} \left( - \left[ \nabla \times \vec{E} \right] \cdot \hat{z} - \sigma^{m}_{z} H_{z} - M_{iz} \right) \\
        & = \frac{1}{\mu_{z}} \left( \frac{\partial E_{x}}{\partial y} - \frac{\partial E_{y}}{\partial x} - \sigma^{m}_{z} H_{z} - M_{iz} \right)
    \end{aligned}
\end{equation}



\subsection{Discretization}%
\subsubsection{Yee locations and leapfrog timing}
On a uniform Cartesian grid, the storage indices $(i,j,k)$ refer to
component-dependent physical locations:
\begin{center}
\renewcommand{\arraystretch}{1.2}
\begin{tabular}{ccc}
\toprule
Component & Spatial location & Sample time \\
\midrule
$E_x$ & $((i+\frac12)\Delta x,j\Delta y,k\Delta z)$ & $n\Delta t$ \\
$E_y$ & $(i\Delta x,(j+\frac12)\Delta y,k\Delta z)$ & $n\Delta t$ \\
$E_z$ & $(i\Delta x,j\Delta y,(k+\frac12)\Delta z)$ & $n\Delta t$ \\
$H_x$ & $(i\Delta x,(j+\frac12)\Delta y,(k+\frac12)\Delta z)$ & $(n+\frac12)\Delta t$ \\
$H_y$ & $((i+\frac12)\Delta x,j\Delta y,(k+\frac12)\Delta z)$ & $(n+\frac12)\Delta t$ \\
$H_z$ & $((i+\frac12)\Delta x,(j+\frac12)\Delta y,k\Delta z)$ & $(n+\frac12)\Delta t$ \\
\bottomrule
\end{tabular}
\end{center}
Electric material parameters are evaluated at electric edges, and magnetic
parameters at magnetic faces. Components sharing $(i,j,k)$ therefore need
not share a physical location or the same effective material value.
The operator $\nabla_h\times$ denotes the discrete curl on this staggered grid.

At each step, first advance $H^{n-\frac12}$ to $H^{n+\frac12}$ using $E^n$
and $M_i^n$, then advance $E^n$ to $E^{n+1}$ using $H^{n+\frac12}$ and
$J_i^{n+\frac12}$. Midpoint fields in the conductivity terms are approximated
by centered averages. The following equalities define the resulting
discrete scheme.

\subsubsection{Electric field}

\begin{equation}
	\begin{aligned}
		\frac{E^{n+1}_{x}(i,j,k) - E_{x}^{n}(i,j,k)}{\Delta t} & =
		\frac{\left[ \nabla_h \times \vec{H} \right]^{n+\frac{1}{2}}(i,j,k) \cdot \hat{x}}{\epsilon_{{x}}(i,j,k)}                                   \\
		                                                       & - \frac{\sigma^{e}_{x}(i,j,k)}{\epsilon_{x}(i,j,k)} E^{n+\frac{1}{2}}_{x}(i,j,k)
		- \frac{1}{\epsilon_{x}(i,j,k)} J^{n+\frac{1}{2}}_{ix}(i,j,k)
	\end{aligned}
\end{equation}

\begin{equation}
	E^{n+\frac{1}{2}}_{x}(i,j,k) \approx \frac{E^{n+1}_{x}(i,j,k)+E^{n}_{x}(i,j,k)}{2}
\end{equation}

\begin{equation}
	\begin{aligned}
		\frac{E^{n+1}_{x}(i,j,k) - E^{n}_{x}(i,j,k)}{\Delta t} & =
		\frac{\left[ \nabla_h \times \vec{H} \right]^{n+\frac{1}{2}}(i,j,k) \cdot \hat{x}}{\epsilon_{{x}}(i,j,k)}                                                     \\
		                                                       & - \frac{\sigma^{e}_{x}(i,j,k)}{\epsilon_{x}(i,j,k)}  \frac{E^{n+1}_{x}(i,j,k)+E^{n}_{x}(i,j,k)}{2}
		- \frac{1}{\epsilon_{x}(i,j,k)} J^{n+\frac{1}{2}}_{ix}(i,j,k)
	\end{aligned}
\end{equation}

\begin{equation}
	\begin{aligned}
		E^{n+1}_{x}(i,j,k) - E^{n}_{x}(i,j,k) & =
		\frac{\Delta t \left[ \nabla_h \times \vec{H} \right]^{n+\frac{1}{2}}(i,j,k) \cdot \hat{x}}{\epsilon_{{x}}(i,j,k)}                                                                                        \\
		                                      & - \frac{\Delta t\sigma^{e}_{x}(i,j,k)}{2\epsilon_{x}(i,j,k)}  E^{n+1}_{x}(i,j,k)- \frac{\Delta t\sigma^{e}_{x}(i,j,k)}{2\epsilon_{x}(i,j,k)}   E^{n}_{x}(i,j,k) \\
		                                      & - \frac{\Delta t}{\epsilon_{x}(i,j,k)} J^{n+\frac{1}{2}}_{ix}(i,j,k)
	\end{aligned}
\end{equation}



\begin{equation}
	\begin{aligned}
		\left(1 + \frac{\Delta t \sigma^{e}_{x}(i,j,k)}{2\epsilon_{x}(i,j,k)} \right) E^{n+1}_{x}(i,j,k) & =
		\frac{\Delta t \left[ \nabla_h \times \vec{H} \right]^{n+\frac{1}{2}}(i,j,k) \cdot \hat{x}}{\epsilon_{{x}}(i,j,k)}                                                                                       \\
		                                                                                                 & + \left( 1 - \frac{\Delta t \sigma^{e}_{x}(i,j,k)}{2\epsilon_{x}(i,j,k)} \right)   E^{n}_{x}(i,j,k) \\
		                                                                                                 & - \frac{\Delta t}{\epsilon_{x}(i,j,k)} J^{n+\frac{1}{2}}_{ix}(i,j,k)
	\end{aligned}
\end{equation}

\begin{equation}
	\boxed{
		f_{e,\alpha}(i,j,k) = \frac{\Delta t \sigma^{e}_{\alpha}(i,j,k)}{2\epsilon_{\alpha}(i,j,k)},
        \qquad \alpha\in\{x,y,z\}.
	}
\end{equation}
The spatial arguments of $f_{e,\alpha}$ are suppressed in the updates below.


\begin{equation}
	\boxed{
		\begin{aligned}
			\left[\nabla_h \times \vec{H} \right]^{n+\frac{1}{2}}(i,j,k) =
			+\left[\frac{H^{n+\frac{1}{2}}_{z}(i,j,k)-H^{n+\frac{1}{2}}_{z}(i,j-1,k)}{\Delta y}
			-\frac{H^{n+\frac{1}{2}}_{y}(i,j,k)-H^{n+\frac{1}{2}}_{y}(i,j,k-1)}{\Delta z} \right] \hat{x} \\
			+\left[\frac{H^{n+\frac{1}{2}}_{x}(i,j,k)-H^{n+\frac{1}{2}}_{x}(i,j,k-1)}{\Delta z}
			-\frac{H^{n+\frac{1}{2}}_{z}(i,j,k)-H^{n+\frac{1}{2}}_{z}(i-1,j,k)}{\Delta x} \right] \hat{y} \\
			+\left[\frac{H^{n+\frac{1}{2}}_{y}(i,j,k)-H^{n+\frac{1}{2}}_{y}(i-1,j,k)}{\Delta x}
			-\frac{H^{n+\frac{1}{2}}_{x}(i,j,k)-H^{n+\frac{1}{2}}_{x}(i,j-1,k)}{\Delta y} \right] \hat{z} \\
		\end{aligned}
	}
\end{equation}

\begin{equation}
    \label{eq:yee-e-x}
	\boxed{
		\begin{aligned}
			E^{n+1}_{x}(i,j,k) & = \frac{1-f_{e,x}}{1+f_{e,x}}  E^{n}_{x}(i,j,k)                                                                                 \\
			                   & + \frac{\Delta t}{(1+f_{e,x})\epsilon_{x}(i,j,k)} \left[ \nabla_h \times \vec{H} \right]^{n+\frac{1}{2}}(i,j,k) \cdot \hat{x} \\
			                   & - \frac{\Delta t}{(1+f_{e,x})\epsilon_{x}(i,j,k)} J^{n+\frac{1}{2}}_{ix}(i,j,k)
		\end{aligned}
	}
\end{equation}

\begin{equation}
    \label{eq:yee-e-y}
	\boxed{
		\begin{aligned}
			E^{n+1}_{y}(i,j,k) & = \frac{1-f_{e,y}}{1+f_{e,y}}  E^{n}_{y}(i,j,k)                                                                                 \\
			                   & + \frac{\Delta t}{(1+f_{e,y})\epsilon_{y}(i,j,k)} \left[ \nabla_h \times \vec{H} \right]^{n+\frac{1}{2}}(i,j,k) \cdot \hat{y} \\
			                   & - \frac{\Delta t}{(1+f_{e,y})\epsilon_{y}(i,j,k)} J^{n+\frac{1}{2}}_{iy}(i,j,k)
		\end{aligned}
	}
\end{equation}

\begin{equation}
    \label{eq:yee-e-z}
	\boxed{
		\begin{aligned}
			E^{n+1}_{z}(i,j,k) & = \frac{1-f_{e,z}}{1+f_{e,z}}  E^{n}_{z}(i,j,k)                                                                                 \\
			                   & + \frac{\Delta t}{(1+f_{e,z})\epsilon_{z}(i,j,k)} \left[ \nabla_h \times \vec{H} \right]^{n+\frac{1}{2}}(i,j,k) \cdot \hat{z} \\
			                   & - \frac{\Delta t}{(1+f_{e,z})\epsilon_{z}(i,j,k)} J^{n+\frac{1}{2}}_{iz}(i,j,k)
		\end{aligned}
	}
\end{equation}



\subsubsection{Magnetic field}

\begin{equation}
	\begin{aligned}
		\frac{H^{n+\frac{1}{2}}_{x}(i,j,k) - H_{x}^{n-\frac{1}{2}}(i,j,k)}{\Delta t} & =
		\frac{\left[ -\nabla_h\times \vec{E} \right]^{n}(i,j,k) \cdot \hat{x}}{\mu_{{x}}(i,j,k)}                                                         \\
		                                                                             & - \frac{\sigma^{m}_{x}(i,j,k)}{\mu_{x}(i,j,k)} H^{n}_{x}(i,j,k)
		- \frac{1}{\mu_{x}(i,j,k)} M^{n}_{ix}(i,j,k)
	\end{aligned}
\end{equation}

\begin{equation}
	H^{n}_{x}(i,j,k) \approx \frac{H^{n+\frac{1}{2}}_{x}(i,j,k)+H^{n-\frac{1}{2}}_{x}(i,j,k)}{2}
\end{equation}

\begin{equation}
	\begin{aligned}
		\frac{H^{n+\frac{1}{2}}_{x}(i,j,k) - H^{n-\frac{1}{2}}_{x}(i,j,k)}{\Delta t} & =
		\frac{\left[ -\nabla_h\times \vec{E} \right]^{n}(i,j,k) \cdot \hat{x}}{\mu_{{x}}(i,j,k)}                                                                                                             \\
		                                                                             & - \frac{\sigma^{m}_{x}(i,j,k)}{\mu_{x}(i,j,k)}  \frac{H^{n+\frac{1}{2}}_{x}(i,j,k)+H^{n-\frac{1}{2}}_{x}(i,j,k)}{2}
		- \frac{1}{\mu_{x}(i,j,k)} M^{n}_{ix}(i,j,k)
	\end{aligned}
\end{equation}

\begin{equation}
	\begin{aligned}
		H^{n+\frac{1}{2}}_{x}(i,j,k) - H^{n-\frac{1}{2}}_{x}(i,j,k) & =
		\frac{\Delta t \left[ -\nabla_h\times \vec{E} \right]^{n}(i,j,k) \cdot \hat{x}}{\mu_{{x}}(i,j,k)}                                                                                                                                           \\
		                                                            & - \frac{\Delta t\sigma^{m}_{x}(i,j,k)}{2\mu_{x}(i,j,k)}  H^{n+\frac{1}{2}}_{x}(i,j,k)- \frac{\Delta t\sigma^{m}_{x}(i,j,k)}{2\mu_{x}(i,j,k)}   H^{n-\frac{1}{2}}_{x}(i,j,k) \\
		                                                            & - \frac{\Delta t}{\mu_{x}(i,j,k)} M^{n}_{ix}(i,j,k)
	\end{aligned}
\end{equation}



\begin{equation}
	\begin{aligned}
		\left(1 + \frac{\Delta t \sigma^{m}_{x}(i,j,k)}{2\mu_{x}(i,j,k)} \right) H^{n+\frac{1}{2}}_{x}(i,j,k) & =
		\frac{\Delta t \left[ -\nabla_h\times \vec{E} \right]^{n}(i,j,k) \cdot \hat{x}}{\mu_{{x}}(i,j,k)}                                                                                                                    \\
		                                                                                                      & + \left( 1 - \frac{\Delta t \sigma^{m}_{x}(i,j,k)}{2\mu_{x}(i,j,k)} \right)   H^{n-\frac{1}{2}}_{x}(i,j,k) \\
		                                                                                                      & - \frac{\Delta t}{\mu_{x}(i,j,k)} M^{n}_{ix}(i,j,k)
	\end{aligned}
\end{equation}

\begin{equation}
	\boxed{
		f_{m,\alpha}(i,j,k) = \frac{\Delta t \sigma^{m}_{\alpha}(i,j,k)}{2\mu_{\alpha}(i,j,k)},
        \qquad \alpha\in\{x,y,z\}.
	}
\end{equation}
The spatial arguments of $f_{m,\alpha}$ are suppressed in the updates below.


\begin{equation}
	\boxed{
		\begin{aligned}
			\left[\nabla_h \times \vec{E} \right]^{n}(i,j,k) =
			\left[\frac{E^{n}_{z}(i,j+1,k)-E^{n}_{z}(i,j,k)}{\Delta y}
			-\frac{E^{n}_{y}(i,j,k+1)-E^{n}_{y}(i,j,k)}{\Delta z} \right] \hat{x} \\
			+\left[\frac{E^{n}_{x}(i,j,k+1)-E^{n}_{x}(i,j,k)}{\Delta z}
			-\frac{E^{n}_{z}(i+1,j,k)-E^{n}_{z}(i,j,k)}{\Delta x} \right] \hat{y} \\
			+\left[\frac{E^{n}_{y}(i+1,j,k)-E^{n}_{y}(i,j,k)}{\Delta x}
			-\frac{E^{n}_{x}(i,j+1,k)-E^{n}_{x}(i,j,k)}{\Delta y} \right] \hat{z} \\
		\end{aligned}
	}
\end{equation}


\begin{equation}
	\boxed{
		\begin{aligned}
			H^{n+\frac{1}{2}}_{x}(i,j,k) & = \frac{1-f_{m,x}}{1+f_{m,x}}  H^{n-\frac{1}{2}}_{x}(i,j,k)                                                    \\
			                             & + \frac{\Delta t}{(1+f_{m,x})\mu_{x}(i,j,k)} \left[ -\nabla_h\times \vec{E} \right]^{n}(i,j,k) \cdot \hat{x} \\
			                             & - \frac{\Delta t}{(1+f_{m,x})\mu_{x}(i,j,k)} M^{n}_{ix}(i,j,k)
		\end{aligned}
	}
\end{equation}

\begin{equation}
	\boxed{
		\begin{aligned}
			H^{n+\frac{1}{2}}_{y}(i,j,k) & = \frac{1-f_{m,y}}{1+f_{m,y}}  H^{n-\frac{1}{2}}_{y}(i,j,k)                                                    \\
			                             & + \frac{\Delta t}{(1+f_{m,y})\mu_{y}(i,j,k)} \left[ -\nabla_h\times \vec{E} \right]^{n}(i,j,k) \cdot \hat{y} \\
			                             & - \frac{\Delta t}{(1+f_{m,y})\mu_{y}(i,j,k)} M^{n}_{iy}(i,j,k)
		\end{aligned}
	}
\end{equation}

\begin{equation}
	\boxed{
		\begin{aligned}
			H^{n+\frac{1}{2}}_{z}(i,j,k) & = \frac{1-f_{m,z}}{1+f_{m,z}}  H^{n-\frac{1}{2}}_{z}(i,j,k)                                                    \\
			                             & + \frac{\Delta t}{(1+f_{m,z})\mu_{z}(i,j,k)} \left[ -\nabla_h\times \vec{E} \right]^{n}(i,j,k) \cdot \hat{z} \\
			                             & - \frac{\Delta t}{(1+f_{m,z})\mu_{z}(i,j,k)} M^{n}_{iz}(i,j,k)
		\end{aligned}
	}
\end{equation}

\section{Time-step stability}
For the standard lossless Yee scheme in a homogeneous isotropic medium,
the Courant--Friedrichs--Lewy (CFL) condition on a uniform Cartesian grid is
\begin{equation}
    \Delta t\leq\frac{1}{v\sqrt{\Delta x^{-2}+\Delta y^{-2}+\Delta z^{-2}}},
    \qquad v=\frac{1}{\sqrt{\epsilon\mu}}.
\end{equation}
Only active spatial directions enter the sum for a reduced-dimensional
simulation. For spatially varying positive isotropic parameters, a
conservative bound uses $v_{\mathrm{bound}}=1/\sqrt{\epsilon_{\min}\mu_{\min}}$,
with independent minima of the discretized absolute material parameters.
In practice choose a safety factor $0<S<1$ and multiply the CFL limit by $S$.
For a cubic three-dimensional vacuum grid with spacing $\Delta$,
\begin{equation}
    \Delta t=S\frac{\Delta}{c_0\sqrt{3}}.
\end{equation}
For passive conductivity, $f\geq0$ gives
$\lvert(1-f)/(1+f)\rvert\leq1$, so the centered loss term is nonamplifying.
This does not remove the propagation CFL restriction. More general material,
boundary, and additional dynamical models require their own stability analysis.

\section{Lumped elements}
\label{sec:lumped}
\subsection{Resistor}
A resistor contributes an electric current to Ampere's law. We first derive
that current for one $x$-directed Yee edge, discretize the resulting Maxwell
equation, and then recover the exact coefficient changes made by the code.
The same derivation applies to $y$- and $z$-directed resistors by permuting
coordinate labels.

We keep the notation of Eq.~\eqref{eq:yee-e-x}: $E_x$ is the electric
field, $\epsilon_x$ is the absolute permittivity, and $\sigma_x^e$ is the
material conductivity at that edge. Only the spatial indices $(i,j,k)$ are
suppressed. A prime on a coefficient means \emph{after adding this resistor};
it does not denote a time derivative. No other impressed source is needed
for this derivation.

\subsubsection{From the resistor current to Ampere's law}
Consider a single field edge of length $\Delta x$. Let $R_{\mathrm{edge}}>0$
be the resistance assigned to that edge. The voltage drop along $+x$ is the
line integral of the electric field. With the Yee edge value representing
the field along that edge,
\begin{equation}
    V_{\mathrm{edge}}=\Delta x\,E_x.
\end{equation}
Ohm's law then gives the current along $+x$:
\begin{equation}
    I_{\mathrm{edge}}=\frac{V_{\mathrm{edge}}}{R_{\mathrm{edge}}}
        =\frac{\Delta x}{R_{\mathrm{edge}}}E_x.
\end{equation}
Ampere's law uses current \emph{density}, not current. This current crosses
the dual face of area $\Delta y\,\Delta z$ associated with the edge, so
\begin{equation}
    \boxed{J_{Rx}=\frac{I_{\mathrm{edge}}}{\Delta y\,\Delta z}
        =\frac{\Delta x}{R_{\mathrm{edge}}\Delta y\,\Delta z}E_x
        =\sigma_{Rx}E_x,\qquad
    \sigma_{Rx}=\frac{\Delta x}{R_{\mathrm{edge}}\Delta y\,\Delta z}.}
    \label{eq:resistor-current}
\end{equation}
The quantity $\sigma_{Rx}$ has units $\mathrm{S/m}$. It is the resistor's
\emph{additional} conductivity at this edge, distinct from the background
material conductivity $\sigma_x^e$.

The $x$ component of Ampere's law, including this resistor current, is
\begin{equation}
    \epsilon_x\frac{\partial E_x}{\partial t}
        =(\nabla\times\vec H)_x-\sigma_x^e E_x-J_{Rx}.
\end{equation}
Substituting Eq.~\eqref{eq:resistor-current} gives
\begin{equation}
    \boxed{\epsilon_x\frac{\partial E_x}{\partial t}
        +(\sigma_x^e+\sigma_{Rx})E_x=(\nabla\times\vec H)_x.}
    \label{eq:resistor-maxwell}
\end{equation}
Thus the resistor enters Maxwell's equation as an additional conduction
current. The permittivity stays the same; the conductivity in the electric
update becomes $\sigma_x^e+\sigma_{Rx}$. No additional term is required in
Faraday's law.

\subsubsection{Discretize the loaded Maxwell equation}
As in the earlier Yee derivation, $E_x$ is stored at integer times and
$\vec H$ at half-integer times. Evaluate Eq.~\eqref{eq:resistor-maxwell}
at $(n+\tfrac12)\Delta t$. Approximate the time derivative by a difference
and the conduction field by the average of its two adjacent samples:
\begin{equation}
    \epsilon_x\frac{E_x^{n+1}-E_x^n}{\Delta t}
    +(\sigma_x^e+\sigma_{Rx})\frac{E_x^{n+1}+E_x^n}{2}
    =(\nabla_h\times\vec H^{n+1/2})_x.
    \label{eq:resistor-discrete}
\end{equation}
The resistor current uses this same average, because
$J_{Rx}^{n+1/2}=\sigma_{Rx}(E_x^{n+1}+E_x^n)/2$.
Multiplying by $\Delta t$ and collecting all $E_x^{n+1}$ terms on the left
gives
\begin{equation}
    \begin{aligned}
    \left[\epsilon_x+(\sigma_x^e+\sigma_{Rx})\frac{\Delta t}{2}\right]
        E_x^{n+1}
    ={}&\left[\epsilon_x-(\sigma_x^e+\sigma_{Rx})\frac{\Delta t}{2}\right]
        E_x^n\\
    &+\Delta t\,(\nabla_h\times\vec H^{n+1/2})_x.
    \end{aligned}
\end{equation}
Dividing by the factor on the left produces the field update:
\begin{equation}
    \boxed{\begin{aligned}
    E_x^{n+1}&=c'_{ee,x}E_x^n
        +c'_{eh,x}(\nabla_h\times\vec H^{n+1/2})_x,\\[2pt]
    c'_{ee,x}&=
        \frac{\epsilon_x-(\sigma_x^e+\sigma_{Rx})\Delta t/2}
             {\epsilon_x+(\sigma_x^e+\sigma_{Rx})\Delta t/2},\\[2pt]
    c'_{eh,x}&=
        \frac{\Delta t}{\epsilon_x+(\sigma_x^e+\sigma_{Rx})\Delta t/2}.
    \end{aligned}}
    \label{eq:resistor-field}
\end{equation}
These are the same two coefficient roles as in Eq.~\eqref{eq:yee-e-x}.
The resistor changes \emph{both} coefficients because its centered current
contains $E_x^{n+1}$ as well as $E_x^n$. Moving its $E_x^{n+1}$ contribution
to the left changes the denominator of the entire update, including the
curl term. Subtracting a current computed only from $E_x^n$ would give a
different time discretization.

\subsubsection{Express the new coefficients using the stored coefficients}
The grid already stores the coefficients before this resistor is attached:
\begin{equation}
    c_{ee,x}=\frac{\epsilon_x-\sigma_x^e\Delta t/2}
                       {\epsilon_x+\sigma_x^e\Delta t/2},\qquad
    c_{eh,x}=\frac{\Delta t}{\epsilon_x+\sigma_x^e\Delta t/2}.
    \label{eq:resistor-old-coefficients}
\end{equation}
They are the $x$ entries of \texttt{c\_ee} and \texttt{c\_eh}.
The helper updates those arrays directly. To see why, divide the numerator
and denominator of $c'_{eh,x}$ in Eq.~\eqref{eq:resistor-field} by the old
denominator $\epsilon_x+\sigma_x^e\Delta t/2$:
\begin{equation}
    \begin{aligned}
    c'_{eh,x}
    &=\frac{\displaystyle
            \frac{\Delta t}{\epsilon_x+\sigma_x^e\Delta t/2}}
           {\displaystyle 1+
            \frac{\sigma_{Rx}\Delta t/2}{\epsilon_x+\sigma_x^e\Delta t/2}}\\[3pt]
    &=\frac{c_{eh,x}}{1+\sigma_{Rx}c_{eh,x}/2}.
    \end{aligned}
\end{equation}
Do the same for $c'_{ee,x}$. Its numerator is the old numerator
$\epsilon_x-\sigma_x^e\Delta t/2$ minus $\sigma_{Rx}\Delta t/2$:
\begin{equation}
    \begin{aligned}
    c'_{ee,x}
    &=\frac{\displaystyle
            \frac{\epsilon_x-\sigma_x^e\Delta t/2}
                 {\epsilon_x+\sigma_x^e\Delta t/2}
            -\frac{\sigma_{Rx}\Delta t/2}
                  {\epsilon_x+\sigma_x^e\Delta t/2}}
           {\displaystyle 1+
            \frac{\sigma_{Rx}\Delta t/2}{\epsilon_x+\sigma_x^e\Delta t/2}}\\[3pt]
    &=\frac{c_{ee,x}-\sigma_{Rx}c_{eh,x}/2}
            {1+\sigma_{Rx}c_{eh,x}/2}.
    \end{aligned}
\end{equation}
Therefore the resistor can be added without recomputing coefficients from
the material arrays:
\begin{equation}
    \boxed{c'_{ee,x}=
        \frac{c_{ee,x}-\sigma_{Rx}c_{eh,x}/2}
             {1+\sigma_{Rx}c_{eh,x}/2},\qquad
        c'_{eh,x}=\frac{c_{eh,x}}{1+\sigma_{Rx}c_{eh,x}/2}.}
    \label{eq:resistor-coefficient-transform}
\end{equation}
This is an algebraic rearrangement of Eq.~\eqref{eq:resistor-field}, with
no extra approximation. After attachment, $c'_{ee,x}$ and $c'_{eh,x}$ become
the stored coefficients for the next call. Applying the same transformation
again adds the next resistor's conductivity while retaining the first.
Two resistors on the same edge therefore give the total conductivity
$\sigma_x^e+\sigma_{R_1x}+\sigma_{R_2x}$, as parallel branches require.
Only the coefficient arrays change; the material arrays are not overwritten.

\subsubsection{Convert the requested terminal resistance to an edge resistance}
For one selected edge, $R_{\mathrm{edge}}=R$, the resistance requested by
the user. A larger resistor region selects several edges. Along $x$, the
code selects $N_s$ edges in series and $N_p$ transverse paths in parallel.
If every edge has resistance $R_{\mathrm{edge}}$, the equivalent network
resistance is
\begin{equation}
    R=\frac{N_s R_{\mathrm{edge}}}{N_p},\qquad
    \boxed{R_{\mathrm{edge}}=R\frac{N_p}{N_s}.}
    \label{eq:resistor-geometry}
\end{equation}
For snapped integer start and end indices, these counts are
\begin{equation}
    N_s=i_e-i_s,\qquad
    N_p=(j_e-j_s+1)(k_e-k_s+1).
\end{equation}
The axial endpoint is excluded from the selected edge indices; the
transverse endpoints are included. Each selected edge still uses its
own dual-face area $\Delta y\,\Delta z$, rather than the entire bounding-box
area. For example, two edges in series and three parallel paths require
$R_{\mathrm{edge}}=3R/2$, since $2(3R/2)/3=R$.
This scaling defines the discrete resistor network; fields in a full
three-dimensional simulation need not be uniform over its paths.

The relevant lines of
\sourcecode{fdtd/lumped_elements.py\#L163-L171}{Resistor.attach\_to\_grid}
now follow directly from Eqs.~\eqref{eq:resistor-geometry} and
\eqref{eq:resistor-current}:
\begin{lstlisting}[language=Python,basicstyle=\small\ttfamily,columns=fullflexible]
resistance = self.resistance * self.parallel / self.series
self.grid._add_electric_load(
    self.index, conductivity=self.edge_length / (resistance * self.edge_area)
)
\end{lstlisting}
Here \texttt{self.resistance} is $R$, the local variable \texttt{resistance}
is $R_{\mathrm{edge}}$, \texttt{edge\_length} is $\Delta x$, and
\texttt{edge\_area} is $\Delta y\,\Delta z$. Thus the helper receives
\texttt{conductivity} $=\sigma_{Rx}$ at exactly the selected $E_x$ entries.

\subsubsection{Match \texttt{\_add\_electric\_load} line by line}
The implementation of
\sourcecode{fdtd/grid.py\#L490-L498}{Grid.\_add\_electric\_load} is
\begin{lstlisting}[language=Python,basicstyle=\small\ttfamily,columns=fullflexible]
def _add_electric_load(self, index, *, permittivity=0, conductivity=0):
    scale = self.c_eh[index] / self.dt
    loss = conductivity * self.dt / 2
    denominator = 1 + (permittivity + loss) * scale
    self.c_ee[index] = (
        self.c_ee[index] + (permittivity - loss) * scale
    ) / denominator
    self.c_eh[index] /= denominator
\end{lstlisting}
For the resistor call, \texttt{permittivity} $=0$ and
\texttt{conductivity} $=\sigma_{Rx}$. Therefore
\begin{equation}
    \begin{aligned}
    \texttt{scale}&=\frac{c_{eh,x}}{\Delta t},\\
    \texttt{loss}&=\frac{\sigma_{Rx}\Delta t}{2},\\
    \texttt{loss}\,\texttt{scale}
        &=\frac{\sigma_{Rx}\Delta t}{2}\frac{c_{eh,x}}{\Delta t}
          =\frac{\sigma_{Rx}c_{eh,x}}{2},\\
    \texttt{denominator}&=1+\frac{\sigma_{Rx}c_{eh,x}}{2}.
    \end{aligned}
\end{equation}
Substituting these expressions into the final two assignments gives
\begin{equation}
    \begin{aligned}
    \texttt{c\_ee[index]}
       &\gets\frac{c_{ee,x}-\sigma_{Rx}c_{eh,x}/2}
                     {1+\sigma_{Rx}c_{eh,x}/2}=c'_{ee,x},\\[2pt]
    \texttt{c\_eh[index]}
       &\gets\frac{c_{eh,x}}{1+\sigma_{Rx}c_{eh,x}/2}=c'_{eh,x}.
    \end{aligned}
\end{equation}
These are exactly Eq.~\eqref{eq:resistor-coefficient-transform}.
The helper accepts a permittivity increment for other loads, but a resistor
adds only conductivity. The resistor derivation therefore uses its
\texttt{permittivity=0} case throughout.

\subsubsection{Where the resistor actually updates the field}
After attachment, \sourcecode{fdtd/grid.py\#L390-L403}{Grid.update\_E}
uses the already modified arrays:
\begin{lstlisting}[language=Python,basicstyle=\small\ttfamily,columns=fullflexible]
self.E *= self.c_ee
self.E += self.c_eh * curl
\end{lstlisting}
At a loaded $x$ edge these two lines mean
\begin{equation}
    E_x^{n+1}=c'_{ee,x}E_x^n
        +c'_{eh,x}(\nabla_h\times\vec H^{n+1/2})_x,
\end{equation}
which is precisely the Maxwell update in Eq.~\eqref{eq:resistor-field}.
\texttt{Resistor} does not implement a separate field correction:
it inherits the empty
\sourcecode{fdtd/lumped_elements.py\#L72-L73}{LumpedElement.update\_E}.
Its current is already included in the two coefficients, so applying it
again would count the same resistor twice.

As a check, remove the magnetic curl and set the background conductivity
to zero. The update reduces to
\begin{equation}
    \frac{E_x^{n+1}-E_x^n}{\Delta t}
       =-\frac{\sigma_{Rx}}{\epsilon_x}
          \frac{E_x^{n+1}+E_x^n}{2}.
\end{equation}
This is the centered discretization of
$\partial E_x/\partial t=-(\sigma_{Rx}/\epsilon_x)E_x$:
the resistor drains the electric energy stored at the edge. The continuous
decay time is $\epsilon_x/\sigma_{Rx}$, while
$|c'_{ee,x}|\leq1$ for passive material and $R_{\mathrm{edge}}>0$.
If the resistor is removed ($R_{\mathrm{edge}}\to\infty$), then
$\sigma_{Rx}\to0$ and both coefficients return to their previous values.

\subsection{Capacitor}
\label{sec:capacitor}
A capacitor contributes a current proportional to the time derivative of
its voltage. As in the resistor derivation, consider one $x$-directed edge
and suppress its spatial indices. Write $\ell=\Delta x$ and
$A=\Delta y\,\Delta z$ for its length and dual-face area. The same equations
apply to the other directions by permuting coordinates. A prime now means
\emph{after adding this capacitor}.

\subsubsection{From the capacitor current to Ampere's law}
For an ideal edge capacitance $C_{\mathrm{edge}}>0$, the charge, voltage,
and current satisfy
\begin{equation}
    Q=C_{\mathrm{edge}}V_{\mathrm{edge}},\qquad
    V_{\mathrm{edge}}=\ell E_x,\qquad
    I_{\mathrm{edge}}=\frac{dQ}{dt}
        =C_{\mathrm{edge}}\ell\frac{\partial E_x}{\partial t}.
\end{equation}
Divide by the dual-face area to obtain the added current density:
\begin{equation}
    \boxed{J_{Cx}=\frac{C_{\mathrm{edge}}\ell}{A}
                    \frac{\partial E_x}{\partial t}
        =\Delta\epsilon_{Cx}\frac{\partial E_x}{\partial t},\qquad
        \Delta\epsilon_{Cx}=\frac{C_{\mathrm{edge}}\ell}{A}.}
    \label{eq:capacitor-current}
\end{equation}
Here $\Delta\epsilon_{Cx}$ has units $\mathrm{F/m}$ and is an absolute
permittivity increment, not a relative permittivity. Substituting this
current into Ampere's law and moving its derivative to the left gives
\begin{equation}
    \epsilon_x\frac{\partial E_x}{\partial t}
        =(\nabla\times\vec H)_x-\sigma_x^e E_x-J_{Cx},
\end{equation}
\begin{equation}
    \boxed{(\epsilon_x+\Delta\epsilon_{Cx})
        \frac{\partial E_x}{\partial t}
        +\sigma_x^e E_x=(\nabla\times\vec H)_x.}
    \label{eq:capacitor-maxwell}
\end{equation}
Thus the capacitor increases electric storage. It adds no conductivity.

\subsubsection{Discretize and collect the field coefficients}
Use the same integer-time electric fields, half-integer-time magnetic
fields, and centered conductivity term as for the resistor:
\begin{equation}
    (\epsilon_x+\Delta\epsilon_{Cx})
        \frac{E_x^{n+1}-E_x^n}{\Delta t}
    +\sigma_x^e\frac{E_x^{n+1}+E_x^n}{2}
        =(\nabla_h\times\vec H^{n+1/2})_x.
\end{equation}
Multiplying by $\Delta t$ and collecting terms yields
\begin{equation}
    \begin{aligned}
    (\epsilon_x+\Delta\epsilon_{Cx}+\sigma_x^e\Delta t/2)E_x^{n+1}
       ={}&(\epsilon_x+\Delta\epsilon_{Cx}-\sigma_x^e\Delta t/2)E_x^n\\
          &+\Delta t\,(\nabla_h\times\vec H^{n+1/2})_x.
    \end{aligned}
\end{equation}
Consequently,
\begin{equation}
    \boxed{\begin{aligned}
    E_x^{n+1}&=c'_{ee,x}E_x^n
        +c'_{eh,x}(\nabla_h\times\vec H^{n+1/2})_x,\\
    c'_{ee,x}&=\frac{\epsilon_x+\Delta\epsilon_{Cx}-\sigma_x^e\Delta t/2}
                         {\epsilon_x+\Delta\epsilon_{Cx}+\sigma_x^e\Delta t/2},\\
    c'_{eh,x}&=\frac{\Delta t}
                         {\epsilon_x+\Delta\epsilon_{Cx}+\sigma_x^e\Delta t/2}.
    \end{aligned}}
    \label{eq:capacitor-field}
\end{equation}
Let $D_x=\epsilon_x+\sigma_x^e\Delta t/2$ denote the old denominator.
Equation~\eqref{eq:resistor-old-coefficients} gives
$c_{eh,x}/\Delta t=1/D_x$. Dividing both new numerators and their
denominator by $D_x$ gives the transformation of the stored arrays:
\begin{equation}
    \boxed{\begin{aligned}
    q_C&=\frac{\Delta\epsilon_{Cx}c_{eh,x}}{\Delta t},\\
    c'_{ee,x}&=\frac{c_{ee,x}+q_C}{1+q_C},\qquad
    c'_{eh,x}=\frac{c_{eh,x}}{1+q_C}.
    \end{aligned}}
    \label{eq:capacitor-coefficients}
\end{equation}
The increment appears with a plus sign in both numerator and denominator
of $c'_{ee,x}$ because it multiplies $E_x^{n+1}-E_x^n$.
For lossless material, $c_{ee,x}=c'_{ee,x}=1$; the capacitor then changes
only $c_{eh,x}$. With conductivity present, both coefficients change.

\subsubsection{Terminal geometry and the exact code calculation}
The counts $N_s$ and $N_p$ are the same as in the resistor section and
are computed by
\sourcecode{fdtd/lumped_elements.py\#L75-L99}{LumpedElement.\_set\_region}.
Series capacitances combine reciprocally, so a network of identical
edge capacitors has
\begin{equation}
    C=\frac{N_p C_{\mathrm{edge}}}{N_s},\qquad
    \boxed{C_{\mathrm{edge}}=C\frac{N_s}{N_p}.}
    \label{eq:capacitor-geometry}
\end{equation}
This scaling is the inverse of the resistor scaling. For two edges in
series and three parallel paths, $C_{\mathrm{edge}}=2C/3$ gives terminal
capacitance $3(2C/3)/2=C$. The implementation in
\sourcecode{fdtd/lumped_elements.py\#L220-L228}{Capacitor.attach\_to\_grid}
is therefore
\begin{lstlisting}[language=Python,basicstyle=\small\ttfamily,columns=fullflexible]
capacitance = self.capacitance * self.series / self.parallel
self.grid._add_electric_load(
    self.index, permittivity=capacitance * self.edge_length / self.edge_area
)
\end{lstlisting}
The local \texttt{capacitance} is $C_{\mathrm{edge}}$, and the helper's
\texttt{permittivity} argument is $\Delta\epsilon_{Cx}$.
In the \sourcecode{fdtd/grid.py\#L490-L498}{electric-load helper},
\texttt{conductivity=0} implies
\begin{equation}
    \begin{aligned}
    \texttt{scale}&=c_{eh,x}/\Delta t,\qquad \texttt{loss}=0,\\
    \texttt{denominator}&=1+\Delta\epsilon_{Cx}\,\texttt{scale}=1+q_C,\\
    \texttt{c\_ee[index]}&\gets(c_{ee,x}+q_C)/(1+q_C),\\
    \texttt{c\_eh[index]}&\gets c_{eh,x}/(1+q_C).
    \end{aligned}
\end{equation}
These are exactly Eq.~\eqref{eq:capacitor-coefficients}. Existing
resistive or capacitive loads are retained because the helper uses the
currently stored coefficients. Like \texttt{Resistor}, \texttt{Capacitor}
inherits the empty \texttt{LumpedElement.update\_E}; the loaded arrays in
\texttt{Grid.update\_E} already implement its entire current contribution.

As a check, with no curl or conductivity, an imposed current density
$J_{ix}$ gives
\begin{equation}
    (\epsilon_x+\Delta\epsilon_{Cx})
        \frac{E_x^{n+1}-E_x^n}{\Delta t}=-J_{ix}^{n+1/2}.
\end{equation}
For one edge, the total storage capacitance is
$C_{\mathrm{edge}}+\epsilon_x A/\ell$. Thus the field retains the
background displacement current in addition to the requested capacitor.
The charging-rate check appears in
\sourcecode{tests/test_physics.py\#L75-L91}{capacitor charging test}.

\subsection{Inductor}
\label{sec:inductor}
An inductor needs a current state: its current cannot be recovered from
the instantaneous field alone. The implemented branch consists of an
inductance $L_{\mathrm{edge}}>0$ and a series resistance
$R_{\mathrm{edge}}\geq0$. Keep $\ell=\Delta x$, $A=\Delta y\,\Delta z$,
and the same voltage and current orientation as above.

\subsubsection{Couple the circuit equation to Maxwell's equation}
The series $RL$ voltage law is
\begin{equation}
    V_{\mathrm{edge}}=L_{\mathrm{edge}}\frac{dI_{\mathrm{edge}}}{dt}
                         +R_{\mathrm{edge}}I_{\mathrm{edge}}.
\end{equation}
Using $V_{\mathrm{edge}}=\ell E_x$ and $I_{\mathrm{edge}}=A J_{Lx}$ gives
\begin{equation}
    \boxed{L_{\mathrm{edge}}\frac{\partial J_{Lx}}{\partial t}
                +R_{\mathrm{edge}}J_{Lx}=\frac{\ell}{A}E_x.}
    \label{eq:inductor-current}
\end{equation}
Ampere's law must include this additional branch current:
\begin{equation}
    \epsilon_x\frac{\partial E_x}{\partial t}+\sigma_x^e E_x
        =(\nabla\times\vec H)_x-J_{Lx}.
    \label{eq:inductor-maxwell}
\end{equation}
Here $\epsilon_x$ and $\sigma_x^e$ may include any static capacitive and
resistive loads already incorporated into the electric coefficients.

\subsubsection{Stagger the current and derive its recurrence}
Store $J_{Lx}$ at half-integer times, alongside the magnetic field.
At the start of the electric update, the available state is
$J_{Lx}^{n+1/2}$. Centering Eq.~\eqref{eq:inductor-maxwell} gives
\begin{equation}
    \epsilon_x\frac{E_x^{n+1}-E_x^n}{\Delta t}
       +\sigma_x^e\frac{E_x^{n+1}+E_x^n}{2}
       =(\nabla_h\times\vec H^{n+1/2})_x-J_{Lx}^{n+1/2}.
\end{equation}
Collecting $E_x^{n+1}$ gives the same coefficients as
Eq.~\eqref{eq:resistor-old-coefficients}, now with an explicit current:
\begin{equation}
    \boxed{E_x^{n+1}=\widetilde E_x^{n+1}-c_{eh,x}J_{Lx}^{n+1/2},\qquad
    \widetilde E_x^{n+1}=c_{ee,x}E_x^n
                     +c_{eh,x}(\nabla_h\times\vec H^{n+1/2})_x.}
    \label{eq:inductor-field}
\end{equation}
The tilde denotes the grid prediction before this branch correction.
Any impressed electric source is applied to that prediction before
\texttt{Inductor.update\_E} runs.

After correcting the field, advance the branch current to $n+3/2$.
Center Eq.~\eqref{eq:inductor-current} at $(n+1)\Delta t$:
\begin{equation}
    L_{\mathrm{edge}}
        \frac{J_{Lx}^{n+3/2}-J_{Lx}^{n+1/2}}{\Delta t}
    +R_{\mathrm{edge}}
        \frac{J_{Lx}^{n+3/2}+J_{Lx}^{n+1/2}}{2}
        =\frac{\ell}{A}E_x^{n+1}.
\end{equation}
Multiplication by $\Delta t/L_{\mathrm{edge}}$ and collection of the
new current yield
\begin{equation}
    (1+d_L)J_{Lx}^{n+3/2}
       =(1-d_L)J_{Lx}^{n+1/2}
         +\frac{\Delta t\,\ell}{L_{\mathrm{edge}}A}E_x^{n+1},
    \qquad d_L=\frac{\Delta t R_{\mathrm{edge}}}{2L_{\mathrm{edge}}}.
\end{equation}
The two additional coefficients are therefore
\begin{equation}
    \boxed{\begin{aligned}
    J_{Lx}^{n+3/2}&=a_L J_{Lx}^{n+1/2}+b_L E_x^{n+1},\\
    a_L&=\frac{1-d_L}{1+d_L},\qquad
    b_L=\frac{\Delta t\,\ell}{L_{\mathrm{edge}}A(1+d_L)}.
    \end{aligned}}
    \label{eq:inductor-coefficients}
\end{equation}
Here $a_L$ is dimensionless, and $b_LE_x$ has units
$\mathrm{A/m^2}$. The series resistance belongs in this current
recurrence. Adding it as an independent conductivity would instead
create a parallel resistor and describe a different circuit.

\subsubsection{Terminal scaling and coefficient calculation in code}
Identical series $RL$ branches combine as
\begin{equation}
    L=\frac{N_s L_{\mathrm{edge}}}{N_p},\qquad
    R=\frac{N_s R_{\mathrm{edge}}}{N_p},\qquad
    \boxed{L_{\mathrm{edge}}=L\frac{N_p}{N_s},\quad
           R_{\mathrm{edge}}=R\frac{N_p}{N_s}.}
    \label{eq:inductor-geometry}
\end{equation}
Thus \texttt{resistance} is the requested terminal series resistance,
and \texttt{inductance} is the terminal inductance before scaling.
The calculation in
\sourcecode{fdtd/lumped_elements.py\#L276-L292}{Inductor.attach\_to\_grid}
matches Eqs.~\eqref{eq:inductor-geometry} and
\eqref{eq:inductor-coefficients}:
\begin{lstlisting}[language=Python,basicstyle=\small\ttfamily,columns=fullflexible]
scale = self.parallel / self.series
inductance = self.inductance * scale
resistance = self.resistance * scale
damping = self.grid.dt * resistance / (2 * inductance)
self._decay = (1 - damping) / (1 + damping)
self._drive = (
    self.grid.dt
    * self.edge_length
    / (inductance * self.edge_area * (1 + damping))
)
self.current_density = np.zeros_like(self.grid.E[self.index])
\end{lstlisting}
The variables \texttt{damping}, \texttt{\_decay}, and \texttt{\_drive}
are $d_L$, $a_L$, and $b_L$. Unlike the resistor and capacitor,
this attachment does not modify \texttt{c\_ee} or \texttt{c\_eh}.
Instead,
\sourcecode{fdtd/lumped_elements.py\#L294-L299}{Inductor.update\_E}
performs both coupled updates:
\begin{lstlisting}[language=Python,basicstyle=\small\ttfamily,columns=fullflexible]
self.grid.E[self.index] -= self.grid.c_eh[self.index] * self.current_density
self.current_density = (
    self._decay * self.current_density + self._drive * self.grid.E[self.index]
)
\end{lstlisting}
Before the first line, \texttt{current\_density} holds
$J_{Lx}^{n+1/2}$; after the second assignment it holds
$J_{Lx}^{n+3/2}$. Attachment initializes this staggered current to zero.
\sourcecode{fdtd/lumped_elements.py\#L273-L274}{Inductor.reset}
clears the state so that a fresh preparation recreates a zero current.

For $R_{\mathrm{edge}}=0$, $a_L=1$ and
$b_L=\Delta t\,\ell/(L_{\mathrm{edge}}A)$, as the ideal voltage law
requires. With no curl or loss, the background edge capacitance is
$C_g=\epsilon_x A/\ell$. Eliminating the current gives
\begin{equation}
    E_x^{n+2}-2E_x^{n+1}+E_x^n
       =-\frac{\Delta t^2}{L_{\mathrm{edge}}C_g}E_x^{n+1}.
\end{equation}
Its oscillatory regime requires
$\Delta t<2\sqrt{L_{\mathrm{edge}}C_g}$, in addition to the propagation
CFL restriction. If a capacitor is present, its added storage belongs
in $C_g$. The discrete oscillation is checked by
\sourcecode{tests/test_physics.py\#L94-L105}{inductor oscillation test};
the current recurrence and reset are checked by
\sourcecode{tests/test_lumped_elements.py\#L43-L56}{inductor recurrence test}.

\subsection{Diode}
\label{sec:diode}
The implemented diode is a Shockley junction with a nonnegative series
resistance $R_s$. Its nonlinear current requires an implicit solve at
each electric update. It occupies exactly one field edge:
$N_s=N_p=1$, so there is no series/parallel rescaling of its parameters.
As above, $\ell=\Delta x$ and $A=\Delta y\,\Delta z$ for an
$x$-directed device.

\subsubsection{Junction law, orientation, and coupling to the field}
Let $p\in\{-1,1\}$ be \texttt{polarity}. Define device voltage and
current positive in the forward direction:
\begin{equation}
    v=p\ell E_x,\qquad J_{Dx}=\frac{pI_D}{A}.
\end{equation}
For junction voltage $v_j$, saturation current $I_s>0$, temperature
$T>0$, and ideality factor $\eta>0$, the model is
\begin{equation}
    \boxed{I_D=I_s\left[\exp\left(\frac{v_j}{V_*}\right)-1\right],
       \qquad v=v_j+R_s I_D,\qquad V_*=\frac{\eta k_B T}{q_e}.}
    \label{eq:diode-law}
\end{equation}
The code stores the ideality-weighted thermal voltage $V_*$ as
\texttt{\_thermal\_voltage}, using
$k_B=1.380649\times10^{-23}\,\mathrm{J/K}$ and
$q_e=1.602176634\times10^{-19}\,\mathrm{C}$.
For $T=300\,\mathrm K$ and $\eta=1$, $V_*\simeq25.852\,\mathrm{mV}$.

Ampere's law contains $-J_{Dx}$, just as it contains the inductor
current. After the ordinary grid prediction $\widetilde E_x^{n+1}$,
including static loads and impressed sources, its centered update is
\begin{equation}
    E_x^{n+1}=\widetilde E_x^{n+1}
                   -c_{eh,x}\frac{pI_D^{n+1/2}}{A}.
\end{equation}
Multiplying by $p\ell$ gives
\begin{equation}
    \boxed{v^{n+1}=v_{\mathrm{pred}}-\gamma I_D^{n+1/2},\qquad
       v_{\mathrm{pred}}=p\ell\widetilde E_x^{n+1},\qquad
       \gamma=\frac{c_{eh,x}\ell}{A}.}
    \label{eq:diode-field}
\end{equation}
The coupling $\gamma$ has units $\Omega$. It uses the final stored
$c_{eh,x}$, so capacitive and resistive loads on the same edge affect
the diode solve. In terms of their effective static parameters,
\begin{equation}
    \gamma=\frac{\Delta t\,\ell}
                    {A(\epsilon_x+\sigma_x^e\Delta t/2)}.
\end{equation}
No constant pair of modified field coefficients can represent the
exponential current for every voltage.

\subsubsection{Develop the scalar implicit equation}
Evaluate the circuit law at the midpoint of the old and corrected
terminal voltages:
\begin{equation}
    \frac{v^n+v^{n+1}}{2}=v_j+R_s I_D^{n+1/2}.
\end{equation}
Substitute Eq.~\eqref{eq:diode-field} and define the midpoint prediction
$m=(v^n+v_{\mathrm{pred}})/2$. Then
\begin{equation}
    m-\frac{\gamma}{2}I_D^{n+1/2}=v_j+R_s I_D^{n+1/2}.
\end{equation}
Substituting the Shockley current produces a single equation for $v_j$:
\begin{equation}
    \boxed{v_j+(R_s+\gamma/2)I_s
                  \left[\exp\left(\frac{v_j}{V_*}\right)-1\right]=m.}
    \label{eq:diode-implicit}
\end{equation}
The term $\gamma/2$ follows from the midpoint field correction; it is
not an extra physical series resistance. Define dimensionless variables
\begin{equation}
    x=\frac{v_j}{V_*},\qquad
    a=\frac{(R_s+\gamma/2)I_s}{V_*},\qquad b=\frac{m}{V_*}.
\end{equation}
The nonlinear residual and its derivative are
\begin{equation}
    \boxed{f(x)=x+a\operatorname{expm1}(x)-b=0,\qquad
        f'(x)=1+a\exp(x)>0.}
    \label{eq:diode-residual}
\end{equation}
For positive coupling, the root is unique and lies in
$[\min(b,0),\max(b,0)]$. Indeed, for $b>0$, $f(0)=-b$ and
$f(b)=a\operatorname{expm1}(b)>0$; for $b<0$ these endpoint signs
reverse. The zero-drive root is $x=0$.

\subsubsection{Match the solver and its coefficients to code}
\sourcecode{fdtd/lumped_elements.py\#L337-L357}{Diode.attach\_to\_grid}
validates the parameters and single-edge geometry, then calculates
\begin{lstlisting}[language=Python,basicstyle=\small\ttfamily,columns=fullflexible]
self._thermal_voltage = (
    self.ideality * 1.380649e-23 * self.temperature / 1.602176634e-19
)
self._gamma = (
    float(self.grid.c_eh[self.index].item()) * self.edge_length / self.edge_area
)
\end{lstlisting}
These are $V_*$ and $\gamma$. Diode attachment leaves the grid
coefficients unchanged. In
\sourcecode{fdtd/lumped_elements.py\#L359-L378}{Diode.\_junction\_voltage},
the argument \texttt{midpoint} is $m$, \texttt{vt} is $V_*$, and
\texttt{a}, \texttt{b}, \texttt{x} are the quantities in
Eq.~\eqref{eq:diode-residual}. The core calculation is
\begin{lstlisting}[language=Python,basicstyle=\small\ttfamily,columns=fullflexible]
a = (self.resistance + self._gamma / 2) * self.saturation_current / vt
b = midpoint / vt
lower, upper = min(b, 0.0), max(b, 0.0)
x = min(b, np.log1p(b / a)) if b > 0 else b
\end{lstlisting}
Each iteration calculates $f(x)$ and $f'(x)$ and proposes the Newton step
\begin{equation}
    x_{\mathrm{new}}=x-\frac{f(x)}{f'(x)}.
\end{equation}
The sign of $f(x)$ first tightens the bracket. A finite Newton candidate
strictly inside the bracket is accepted; otherwise the code bisects it.
The implementation uses \texttt{expm1} and \texttt{log1p} to retain
accuracy near zero and caps exponential arguments at $700$ to limit
overflow. That cap is a numerical safeguard; Eq.~\eqref{eq:diode-law}
states the intended junction model. In numerical voltage values, the
stopping condition is
\begin{equation}
    |f(x)|V_*\leq\texttt{tolerance}\,(1+|m|).
\end{equation}
The default \texttt{tolerance} is $10^{-12}$ and the iteration limit is
$100$. A failed solve raises \texttt{RuntimeError}; nonfinite or complex
time-domain fields are rejected.

Once $v_j=xV_*$ is known, Eq.~\eqref{eq:diode-implicit} also provides
the current without evaluating another exponential:
\begin{equation}
    I_D^{n+1/2}=\frac{m-v_j}{R_s+\gamma/2},\qquad
    v^{n+1}=v_{\mathrm{pred}}-\gamma I_D^{n+1/2},\qquad
    E_x^{n+1}=\frac{p v^{n+1}}{\ell}.
\end{equation}
These are the assignments in
\sourcecode{fdtd/lumped_elements.py\#L380-L394}{Diode.update\_E}:
\begin{lstlisting}[language=Python,basicstyle=\small\ttfamily,columns=fullflexible]
predicted = self.polarity * self.edge_length * field.real
midpoint = (self.previous_voltage + predicted) / 2
junction = self._junction_voltage(midpoint)
current = (midpoint - junction) / (self.resistance + self._gamma / 2)
voltage = predicted - self._gamma * current
self.grid.E[self.index] = self.polarity * voltage / self.edge_length
self.previous_voltage = voltage
\end{lstlisting}
At entry, \texttt{previous\_voltage} represents $v^n$; the final
assignment stores $v^{n+1}$ for the next step. Attachment and
\texttt{reset} initialize it to zero; a nonzero initial field does not
automatically initialize this separate voltage history.
The method refreshes \texttt{\_gamma} from \texttt{c\_eh} before each
solve because later attachments may have added parallel static loads.

As a check, set $R_s=0$ and $v^n=0$. Then
$v_j=v^{n+1}/2$, and the update must satisfy
\begin{equation}
    v^{n+1}-v_{\mathrm{pred}}
        +\gamma I_s\operatorname{expm1}
                      \left(\frac{v^{n+1}}{2V_*}\right)=0.
\end{equation}
This is the residual checked for every direction and polarity in
\sourcecode{tests/test_lumped_elements.py\#L63-L82}{diode residual test}.
The model includes neither junction capacitance nor reverse breakdown.
Dynamic elements sharing an edge are applied sequentially by
\sourcecode{fdtd/grid.py\#L390-L403}{Grid.update\_E}; the equations above
describe each element's correction, not a simultaneous nonlinear solve
for all overlapping dynamic branches.

% \begin{figure}[h]
% 	\centering
% 	\begin{subfigure}[b]{0.8\textwidth}
% 		\centering
% 		\def\svgwidth{\columnwidth}
% 		\includesvg[figures/]{image1}
% 		\caption{}
% 	\end{subfigure}

% 	\begin{subfigure}[b]{0.8\textwidth}
% 		\centering
% 		\def\svgwidth{\columnwidth}
% 		\includesvg[figures/]{image1}
% 		\caption{}
% 	\end{subfigure}
% 	\caption{}
% 	\label{}
% \end{figure}


\end{document}

